\documentclass{beamer}
\usepackage{lmodern}

\usetheme{CambridgeUS} % try Madrid, Pittsburgh
\usecolortheme{beaver}
\usefonttheme[onlymath]{serif} % try "professionalfonts"


\graphicspath{{./figure/}}

\setbeamertemplate{itemize items}[default]
\setbeamertemplate{enumerate items}[default]

\usepackage{amsmath, amsfonts, latexsym, mathtools, extarrows}
\usepackage{centernot} % for ``not implies'' symbol
\DeclareMathOperator*{\argmin}{arg\,min}
\DeclareMathOperator*{\argmax}{arg\,max}

\usepackage{pifont}
\usepackage{hyperref}
\usepackage{comment}

\usepackage[normalem]{ulem}
\newcommand{\middlewave}[1]{\raisebox{0.5em}{\uwave{\hspace{#1}}}}

\usepackage{graphicx, subcaption}

\usepackage{tikz}

\usepackage{verbatim}
\usepackage{algorithm}
\usepackage[noend]{algpseudocode}

%hyper ref
\usepackage{hyperref}

% table
\usepackage{amstext} % for \text macro
\usepackage{array}   % for \newcolumntype macro
\newcolumntype{L}{>{$}l<{$}} % math-mode version of "l" column type
\newcolumntype{C}{>{$}c<{$}} % math-mode version of "c" column type
\newcolumntype{R}{>{$}r<{$}} % math-mode version of "r" column type
\newcommand{\pno}[1]{\textcolor{blue}{\scriptsize [Problem: #1]}}
\newcommand{\set}[1]{\{#1\}}
\newcommand{\cmark}{\textcolor{red}{\ding{51}}}
\newcommand{\xmark}{\textcolor{red}{\ding{55}}}
%%%%%%%%%%%%%%%%%%%%%%%%%%%%%%%%%%%%%%%%%%%%%%%%%%%%%%%%%%%%%%
% for fig without caption: #1: width/size; #2: fig file
\newcommand{\fignocaption}[2]{
	\begin{figure}[htp]
		\centering
		\includegraphics[#1]{#2}
	\end{figure}
}
% for fig with caption: #1: width/size; #2: fig file; #3: fig caption
\newcommand{\fig}[3]{
	\begin{figure}[htp]
		\centering
		\includegraphics[#1]{#2}
		\caption[labelInTOC]{#3}
	\end{figure}
}
\newcommand{\titletext}{Programming Assignment Lecture \uppercase\expandafter{\romannumeral1}}

%%%%%%%%%%%%%%%%%%%%
\title{\titletext}
\subtitle{Introduction}
\author{Xie}
\institute{xiemhemail@gmail.com}
\date{Sep 8th, 2017}
\titlegraphic{\includegraphics[height = 2cm]{figure/qrcode.png}}



\AtBeginSection[]{
	\begin{frame}[noframenumbering, plain]
		\frametitle{\titletext}
		\tableofcontents[currentsection, sectionstyle=show/shaded, subsectionstyle=show/show/hide]
	\end{frame}
}
%%%%%%%%%%
\begin{document}
	\maketitle
	
\begin{frame}{QQ Group:278741068}
	\fignocaption{width=200pt}{qqqrcode.jpg}
\end{frame}

\begin{frame}{Basic Info}
	\begin{itemize}
		\item Course Content
		\begin{itemize}
			\item PA1-PA??
		\end{itemize}
		\item Time
		\begin{itemize}
			\item Friday each odd weak
		\end{itemize}
		\item Assessment methods
			\begin{itemize}
				\item Submit project and report before \textbf{deadline}.
				\item An act of plagiarism is absolutely forbidden.
				\item For more details, see \href{https://nju-ics.gitbooks.io/ics2017-programming-assignment/content/submit.html}{\uwave{\uppercase{submit requirement webpage}}}.
			\end{itemize}
		\item \alert{Instruction}
		\begin{itemize}
			\item https://nju-ics.gitbooks.io/ics2017-programming-assignment/
			\item \alert{\huge Please check the "news" module of website at least once each day.}
		\end{itemize}
	\end{itemize}
\end{frame}


\begin{frame}{Teacher and T.A.s}
	\begin{description}
		\item[Teacher] Liang Wang(\href{mailto:wl@nju.edu.cn}{wl@nju.edu.cn})
		\item[T.A.] Ruize Tang(\href{mailto:151220100@smail.nju.edu.cn}{151220100@smail.nju.edu.cn})\\
					Minhui Xie(\href{mailto:xiemhemail@gmail.com}{xiemhemail@gmail.com})
	\end{description}
\end{frame}

\begin{frame}{Task load index}
	Here are 3 questionnaires to assess your perceived workload.\alert{We sincerely hope you finish them} in each stage for our teaching research.But it isn't mandatory,\textbf{all by voluntary}!
	\begin{description}
		\item 	[\href{https://github.com/xmh645214784/ics2017Slides/blob/master/introduction/Q1.doc}{Q1.doc}] Please fill it up before the start of PA \textbf{and} fill it up when submitting \textbf{each} large stage.
		\item	[\href{https://github.com/xmh645214784/ics2017Slides/blob/master/introduction/Q2.doc}{Q2.doc}] id.
		\item	[\href{https://github.com/xmh645214784/ics2017Slides/blob/master/introduction/Q3.doc}{Q3.doc}] Please fill it up when submitting \textbf{each} large stage.
	\end{description}
\begin{alertblock}{Submission format}
	Compress them with your report and PA. Submission of the first time of Q1 and Q2 (\textit{i.e. Before the start of PA}) can be postponed to the submission of whole PA1. But make sure you have distinguish them in file name such like Q1\_1st.doc,Q1\_2nd.doc.
\end{alertblock}
\end{frame}

\begin{frame}{Resources}
	\begin{description}
		\item [Plantform and tools] IA-32 + GNU/Linux + gcc + C
		\item [Guidebook] https://nju-ics.gitbooks.io/ics2017-programming-assignment/
		\item [Skeleton]https://github.com/NJU-ProjectN/ics-pa
	\end{description}
	\begin{alertblock}{Tip}
		You can download the PDF or epub version of guide in github.
	\end{alertblock}
\end{frame}


\section{From Computer System, to ICS, to PA}

\begin{frame}
	\centering{\huge Why we need learn ICS?}
\end{frame}

\begin{frame}[fragile]{Motivation}
\begin{exampleblock}{Question}
\begin{verbatim}
int main()
{
    printf("Hello World");
    return 0;
}
\end{verbatim}
What the computer are doing when you execute the program above？
\end{exampleblock}

\begin{alertblock}{Tip}
	This may appear in exams.
\end{alertblock}

\end{frame}




\begin{frame}{System Stack}
\begin{table}
	\centering
	\begin{tabular}{|c|}
		\hline
		Application \\
		\hline
		Algorithm\\
		\hline
		Programming Language\\
		\hline
		Operating System/Virtual Machines\\
		\hline
		\alert{Instruction Set Architecture}		\\
		\hline
		Micro-architecture\\
		\hline
		Register-Transfer Level\\
		\hline
		Gates\\
		\hline
		Circuits\\
		\hline
		Devices\\
		\hline
		Physics\\
		\hline
	\end{tabular}
\end{table}
\end{frame}



\begin{frame}{What you'll get after finishing PA}
	\begin{description}
		\item[You'll]

		\begin{itemize}
		\item	\textbf{Get Systems thinking}
		\item	Understand how program run on a computer
		\item	Enhance \textbf{coding} ability
		\item	Prepare for later courses (OS,Compiling)
		\end{itemize}
		\pause
		\item[Way]Complete a tiny but entire computer system and run program on it.
		\pause
		\item[PA] \alert{\huge NEMU}(\textit{\uwave i.e. NJU Emulator})
	\end{description}
\begin{exampleblock}{Question}
	\huge What is an \textit{emulator}?
\end{exampleblock}
\end{frame}


\begin{frame}{What is an \textit{emulator}?}
	\begin{block}{Emulator --Wikipedia}
		In computing, an \textbf{emulator} is hardware or software that enables one computer system (called the \textit{host}) to behave like another computer system (called the \textit{guest}).

		An emulator typically enables the host system to run software or use peripheral devices designed for the guest system.
	\end{block}
\end{frame}


\begin{frame}{What is an \textit{emulator}?}
	\fignocaption{scale=0.4}{supermario.jpg}
	\begin{columns}
		\begin{column}{0.4\textwidth}
			\fignocaption{scale=0.5}{fc.jpg}
		\end{column}

		\begin{column}{0.2\textwidth}
			\begin{tikzpicture}[line width=5pt]
				%\draw (1,1) node {FC Emulator}
				\draw[red] [<-] (0,0) --(2,0);
			\end{tikzpicture}
		\end{column}

		\begin{column}{0.4\textwidth}
			\fignocaption{scale=0.5}{pc.jpg}
		\end{column}
	\end{columns}
\end{frame}

\begin{frame}{What is an \textit{emulator}?}
	\fignocaption{scale=0.4}{debian.png}
	\begin{columns}
		\begin{column}{0.4\textwidth}
			\fignocaption{scale=0.4}{pc.jpg}
		\end{column}

		\begin{column}{0.2\textwidth}
			\begin{tikzpicture}[line width=5pt]
			%\draw (1,1) node {FC Emulator}
			\draw[red] [<-] (0,0) --(2,0);
			\end{tikzpicture}
		\end{column}

		\begin{column}{0.4\textwidth}
		%	\fignocaption{scale=0.1}{virtualbox.jpg}
			\fignocaption{scale=0.2}{vmware.jpg}
		\end{column}
	\end{columns}
\end{frame}


\begin{frame}{}
\fignocaption{scale=0.3}{emulate.png}
\begin{exampleblock}{Question}
	Does the object know he is living in a virtualized environment?
\end{exampleblock}

\end{frame}


\begin{frame}{Task of PA}
	\begin{block}{Task of PA}
		To realize \alert{NEMU}, a  simplified x86 system-wide \textbf{emulator}.
	\end{block}
%\begin{columns}
%	\begin{column}{0.30\textwidth}
%		\begin{table}[htbp]
%			\begin{tabular}{|c|}
%				\hline
%				"Hello world" program\\
%				\hline
%				GNU/Linux\\
%				\hline
%				Computer hardware\\
%				\hline
%			\end{tabular}
%		\end{table}
%
%		\begin{table}[htbp]
%		\begin{tabular}{|c|}
%			\hline
%			SuperMario\\
%			\hline
%			Simulated NES hardware\\
%			\hline
%			NES Emulator\\
%			\hline
%			GNU/Linux\\
%			\hline
%			Computer hardware\\
%			\hline
%		\end{tabular}
%	\end{table}
%	\end{column}
%
%	\begin{column}{0.30\textwidth}
%	\begin{table}[htbp]
%		\begin{tabular}{|c|}
%			\hline
%			"Hello world" program\\
%			\hline
%			Micro OS\\
%			\hline
%			Simulated x86 hardware\\
%			\hline
%			NEMU\\
%			\hline
%			GNU/Linux\\
%			\hline
%			Computer hardware\\
%			\hline
%		\end{tabular}
%	\end{table}
%	\end{column}
%
%\end{columns}
\fignocaption{scale=0.3}{emulator.png}
\end{frame}

%\begin{frame}{学生体会}
%	2017.1.11 2:14AM
%
%	不相信所谓的执念, 本已打算放弃, 交换了bt指令的源操作数和目的操作数, 出现了在别人家的电脑上出现过无数次的画面, 耳边是Anan Ryoko的Refrain, 却已经难以抑制地在只有一个人的寝室里失声痛哭.
%
%
%	从汇编语言课程开始听闻PA实验, 举步维艰地进行每一个阶段, 无数个不眠的深夜, 第二天拖着疲惫的身体睁着朦胧的睡眼去上课, 我一次次地怀疑自己, 一次次想要放弃. 现在想来, 这也印证了”成长是一个痛苦的过程”这段文字. 如果说承认自己的软弱, 是成长的第一步. 那么可以恬不知耻地说, 我已经真正从19岁走向20岁了. 是的, 你在专业上的技不如人, 迟早有一天会找上来, 我已经真正感受到了. 既然选择了远方, 便只顾风雨兼程. …当时有很多规划, 现在看来都在不断调整, 然而有一点没有变, 我的人生职业目标: 成为一名IT精英. 仙剑的运行作为IT路上的第一个里程碑吧.
%\end{frame}

\begin{frame}{Transition}
	\centering {\huge{So how we emulate a computer?}}
\end{frame}

\section{How we emulate a computer? $\longrightarrow$ The story of computer}
\begin{frame}{DLC}
	\begin{itemize}
		\item Most of you have completed the Digital Logic Circuit Course.(Taught by zzs whj)
			\begin{itemize}
				\item Adder
				\item Register
				\item Multiplexer
				\item \textit{etc.}
			\end{itemize}
		\pause
		\item Logic Gate $\longrightarrow$ digital logic device $\longrightarrow$ Computer
	\end{itemize}
\end{frame}

\begin{frame}{The simplest computer- Turing machine}
\begin{block}{Ultimate goal of computer}
	Run programs on it.
\end{block}
	\begin{description}
		\centering
		\item[To place programs] Memory
		\item[To process data]Adder
		\item[To store temporary results efficiently]Reg
	\end{description}
	\fignocaption{scale=0.5}{trm.png}
\end{frame}

\begin{frame}{The simplest computer- Turing machine}
	\begin{itemize}
		\item TRM = Reg + Adder + Mem
		\begin{itemize}
			\item These are all what we learned in DLC!
		\end{itemize}
		\item Computers read data from Mem, store them to Reg, do some calculate and write the result to Mem.
	\end{itemize}
	\fignocaption{scale=0.5}{trm.png}
\end{frame}

\begin{frame}{State of TRM}
	\begin{itemize}
		\item Reg and Mem are all Sequential Logic Device.They can store value (i.e. state).
		\item The process of computer working = These sequential logical devices  \alert{transfer one state to another state}
	\end{itemize}
	\fignocaption{scale=0.3}{stateTRM.png}
\end{frame}

\begin{frame}{Instructions of TRM}
	\begin{itemize}
		\item This large state machine does have abundant states.
		\item But how one state transfer to another is not at liberty.
		\item \alert{Instructions = Commands that guide the state transition of computer}
		\begin{itemize}
			\item R/W Mem ,Computing
		\end{itemize}
		\item \alert{Program = A lot of instructions}
	\end{itemize}
\fignocaption{scale=0.4}{stateTRM2.png}
\end{frame}

\begin{frame}{Instruction Set}
	\begin{itemize}
		\item \alert{Instruction Set = All things computer can do}
		\begin{itemize}
			\item	e.g. inc,dec,jne
		\end{itemize}

\end{itemize}

\begin{block}{Question}
	What can TRM do with just these 3 instructions?
\end{block}
\pause
\begin{itemize}
		\item Computability theory shows that TRM can do anything!
		\begin{description}
			\item [loop]  jne
			\item [add] inc,inc,inc,inc,inc...
			\item [multiply] add,add,add,add...
			\item [\textbf{NEMU}] Wow!
		\end{description}

	\end{itemize}

\end{frame}

\begin{frame}{Instruction Set}
	\begin{itemize}
		\item However,the speed of only 3 instructions-version TRM is so slow.
		\item Why do we add more powerful instructions to the computer?
		\item \alert{Modern computer Instructigiton Set}(x86, mips…)
		\begin{itemize}
			\item mul,div, cmp, bit operation, string manipulation…
			\item Adder$\longrightarrow$ALU
		\end{itemize}

	\end{itemize}
	\fignocaption{scale=0.4}{modernisa.png}
\end{frame}

\begin{frame}{Input \& Output}
	\begin{block}{Question}
		Now,it seems that we have a powerful computer.Are there any unsatisfactory drawbacks?
	\end{block}
	\begin{itemize}
		\pause
		\item It can only do computing!
		\item To interact with the outside world, \alert{we need to add I/O devices}
		\item plus with some relative instructions.
	\end{itemize}

	\fignocaption{scale=0.4}{io.png}
\end{frame}

\begin{frame}{Interrupt \& Exception}
	\begin{itemize}
		\item Besides running regular program, computer should deal with emergencies whenever possible.
		\item Such as,
		\begin{itemize}
			\item Inner exception: division by 0, access violation, trap etc.
			\item Outer interrupt:Keyboard, device ready,etc.
		\end{itemize}
		\item	\alert{We need add asynchronous processing unit to deal with emergencies.}
	\end{itemize}
	\fignocaption{scale=0.4}{interrupt.png}
\end{frame}

\begin{frame}{Multitask}
\begin{block}{Question again}
	Now,it seems that we have a powerful computer.Are there any unsatisfactory drawbacks?
\end{block}
\pause
\begin{itemize}
	\item Monotask $\longrightarrow$ Mutitask
	\item \alert{Time-division multiplexing}: Improve the efficiency of resources usages.
	\begin{itemize}
		\item CPU, I/O Device: Take turns!
	\end{itemize}
\end{itemize}
\begin{block}{Question}
	Can tasks use memory exclusively by turns?What else should we do in practice?
\end{block}
\pause
\begin{alertblock}{Answer}
	Certainly not.Exclusive memory using will cause inefficiency.\\
\alert{We should add MMU for memory protection}.
\end{alertblock}

%	\fignocaption{scale=0.4}{am.png}
\end{frame}

\begin{frame}{AM(Abstract Machine)}
\centering{\huge{AM = TRM + IOE + ASYE + PTE + MPE}}
\begin{itemize}
	\item TRM(Turing Machine)
	\item IOE(I/O Extension)
	\item ASYE(Asynchronous Extension)
	\item PTE(Protection Extension)
	\item MPE(Multi-Processor Extension) (not included in PA)
\end{itemize}
\fignocaption{scale=0.4}{am.png}
\end{frame}

%\begin{frame}{Meaning of AM}
%	\begin{itemize}
%		\item Devices vary from computer to computer, and ISA.
%		\item If we predefine some APIs which computers can do, then all computers realize them in bottom,we'll get a unified Abstract Machine.
%		\begin{itemize}
%			\item Then programs can directly run on AM with no care about the differences among ISAs.
%		\end{itemize}
%
%		\item ISA = Interface between software and hardware
%		\item APIs of AM = Software of ISA.
%		\begin{itemize}
%			\item 	So, AM is the software closest to hardware.
%		\end{itemize}
%	\end{itemize}
%
%
%\end{frame}

\begin{frame}{How to realize NEMU?}
	\begin{block}{}
		By now, we have known what a modern computer should look like, then how we emulate it? Or what is the constituent of NEMU?
	\end{block}
\pause
	\begin{block}{What tools we already have?}
		\pause
		\begin{itemize}
			\item Variables,data structure
			\item Algorithm
		\end{itemize}


	\end{block}
	\begin{description}
		\item[Devices] C variables, arrays, link list, etc.
		\item[Logic functions] Algorithmic functions
		\item[Runtime support] Loop
	\end{description}
\end{frame}

\begin{frame}[fragile]{E.g. cpu}
\begin{verbatim}
struct CPU_State
{
    rtlreg_t eax, ecx, edx, ebx, esp, ebp, esi, edi;
    vaddr_t eip;
    EFLAGS;
    CRs;
    ......
}
\end{verbatim}
\fignocaption{scale=0.2}{cpu.jpg}
\end{frame}

\begin{frame}[fragile]{E.g. memory}
\begin{verbatim}
uint8_t pmem[PMEM_SIZE];
\end{verbatim}
	\fignocaption{scale=0.2}{memory.jpg}
\end{frame}

\begin{frame}{The simplest computer- Turing machine}
	\begin{columns}
		\begin{column}{0.5\textwidth}
			\begin{block}{Architecture of the simplest computer }
					\begin{description}
					\item[To place programs] Memory
					\item[To process data]Adder
					\item[To store temporary results efficiently]Reg
				\end{description}
			\end{block}
			\fignocaption{scale=0.4}{trm.png}

		\end{column}
		\begin{column}{0.5\textwidth}
				\begin{block}{Working mode of the simplest computer}
				\begin{itemize}
					\item Fetch instruction from Mem using PC.
					\item Execute instruction.
					\item Update PC.
				\end{itemize}
			\end{block}
			\fignocaption{scale=0.5}{tapemachine.jpg}
	\end{column}
	\end{columns}
\end{frame}


\begin{frame}{TRM in NEMU}
		\begin{block}{Architecture of the simplest computer }
			\begin{description}
				\item[Memory] \textbf{pmem array}(128MB nemu/src/memory/memory.c)
				\item[Adder(ALU)] \textbf{functions}
				\item[Reg]	\textbf{CPU\_state struct} (nemu/include/cpu/reg.h)
			\end{description}
		\end{block}
		\begin{block}{Working mode of the simplest computer}
			\textbf{Function cpu\_exec() }(nemu/src/monitor/cpu-exec.c)
			\begin{itemize}
				\item Fetch instruction from Mem using PC.
				\item Execute instruction.
				\item Update PC.
			\end{itemize}
		\end{block}
\end{frame}

\section{Help you do it! $\longrightarrow$ Brand new PA based on AM}
\begin{frame}[fragile]{Structure of PA}
\begin{verbatim}
    ics2017
    |---nanos-lite    # mini operating system kernel
    |---navy-apps     # apps
    |---nemu          # NEMU
    |---nexus-am      # abstract machine
\end{verbatim}
\end{frame}

\begin{frame}{Structure of PA}
	\fignocaption{scale=0.4}{brandnewPA.png}
\end{frame}

\begin{frame}{PA0 and PA1}
	\begin{itemize}
		\item PA0 - Environment configuration
		\item PA1 - Simple monitor
	\end{itemize}
	\fignocaption{scale=0.5}{trm.png}
\end{frame}

\begin{frame}{PA2-Von Neumann Computer System (TRM+IOE)}
	\begin{itemize}
		\item Instructions system + I/O
		\item Run interesting apps on NEMU
	\end{itemize}
\fignocaption{scale=0.3}{io.png}
\begin{columns}
	\begin{column}{0.2\textwidth}
		\fignocaption{scale=0.35}{typinggame.png}
	\end{column}
\begin{column}{0.2\textwidth}
	\fignocaption{scale=0.35}{supermario.png}
\end{column}
\begin{column}{0.2\textwidth}
	\fignocaption{scale=0.35}{pal.png}
\end{column}
\end{columns}
\end{frame}


\begin{frame}{PA3- Asynchronous Extension}
	\begin{itemize}
		\item Run tiny operating system Nanos-lite on NEMU.
		\begin{itemize}
			\item ramdisk, fs, raw img loader
		\end{itemize}
		\item Implement some library functions which packing the system call .
		\item App calls library functions and run PAL.
	\end{itemize}
	\fignocaption{scale=0.4}{interrupt.png}
\end{frame}

\begin{frame}{PA4- Multitask(Protection Extension)}
	\begin{itemize}
	\item Virtual storage.
	\item Interruption.
	\item Run pal,typing game,clock at the same time on Nanos-lite.
\end{itemize}
\fignocaption{scale=0.4}{am.png}
\end{frame}

\begin{frame}{PA5- Performance}
	\begin{itemize}
		\item Elf32 loader
		\item cache
		\item TLB
		\item profiler, performance optimization
		\item JIT \textit{(i.e just-in-time)} compilation
	\end{itemize}
\end{frame}

\begin{frame}
	\fignocaption{scale=0.4}{estimatedTime.png}
\end{frame}

\begin{frame}{The End}
	\fignocaption{scale=0.8}{studentFeedback.png}
\end{frame}


\end{document}
